% !TEX root = ../thesis-example.tex
%
\chapter{Conclusions}
\label{sec:conclusions}

\section{Limitations}

We weren't able to properly implement a \ac{PC} face detection model for the following reasons:
\begin{itemize}
	\item Training a 3D based model is expensive computation wise and time wise so we could only use pretrained things.
	\item Face detection is not addressed specifically in tasks so we tried to use the pedestrian category.
	\item Wherever PyTorch/TensorFlow did not provide all of the foundational blocks required,
	      only CUDA GPU implementations were provided without a CPU or AMD GPU equivalent.
\end{itemize}

\par The 3D AI-based Object Detection and general \ac{PC} processing python ecosystem is not that mature yet.
There are great tools that are either still at an early research phase or not properly maintained. This makes them unsuitable for software in production environments. A notable exception was the \emph{Open3D} package.

\par Face detection as a specific task is not that popular in Point cloud related settings. Mostly pedestrians in an outdoor setting are actually the target of the latest models. There has been more focus in face recognition tasks as point clouds can provide richer features of a face than the two-dimensional equivalent and can therefore improve the effectiveness of face recognition systems.

\par Most of the attention is on traffic settings, other use cases remain mostly unexplored e.g. indoors 3D cameras.
In \cite{wang_voxel-rcnn-complex_2022} the authors improve upon an existing model to make it more suitable for complex traffic conditions.
\section{Discussion}
Anything in 3D Point-Clouds is more complex than the 2D equivalent for many reasons. The non-grid-like structure makes simple operations more advanced to implement, the sparseness of \ac{PC}s necessitates extra attention when it comes to statistical metrics used and the significant payload size drives up the computation cost considerably.
Being able to efficiently and effectively provide privacy awareness in \ac{PC}s could enable work on traffic settings by reducing the legal process that datasets have to go through, and enable more people-centered applications of \ac{PC} work.

\section{Future Work}

Several research directions occur from the research conducted in this thesis.

\textbf{Evaluation of 2D vs 3D detection models.}
A systematic comparison between 2D and 3D object detection models for face and people
detection in both indoor and outdoor settings would provide more insight into the
trade-offs between the two approaches. Such an evaluation should consider detection
accuracy, robustness to lighting conditions, occlusion handling, and computational cost.

\textbf{Optimization of 3D detection models for embedded systems.}
Investigating how to optimize 3D detection models for deployment on embedded
and low-powered systems, is an important next step. Currently models achieve
inference times around 30-60 FPS but with special server/desktop hardware
as noted in \cite{zimmer_survey_2022}.

\textbf{Color-independent model training.}
Training detection models on colorless point clouds and evaluating whether they
can match the performance of their color-dependent counterparts would be a valuable
contribution. Color information in existing datasets can introduce biases tied to
specific environments or demographics. A model that relies solely on geometric
features would generalize better across different sensors and scenarios, and would
be applicable in settings where color data is unavailable or unreliable.

\textbf{End-to-end 3D video anonymization pipelines.}
Building on work such as \cite{Zhao_Zhang_Demiris_2024}, developing end-to-end
pipelines for 3D video anonymization would enable privacy-compliant data collection
at scale. Such pipelines would need to handle temporal and visual consistency across frames,
real-time performance, and robustness to challenging conditions, representing a
significant but impactful engineering and research challenge.
